\documentclass[10pt,letterpaper]{article}
\usepackage{cornell-notes}

\title{Clause 03.11: Correspondence}
\author{}
\date{}
\setDocTitle{Clause 03.11: Correspondence}
\setDocSubtitle{ISO/IEC/IEEE 42010:2022}
\setDocOwner{ISO 42010 Cornell Notes}
\setDocDate{\today}
\setCornellCollection{ISO/IEC/IEEE 42010:2022 Cornell Notes}
\setCornellUnitType{Clause}
\setCornellUnitNumber{03.11}
\setCornellUnitTitle{Correspondence}
\hypersetup{pdftitle={Clause 03.11: Correspondence},pdfsubject={Cornell notes},pdfauthor={}}

\begin{document}
\maketitle
\begin{CornellOverviewBox}{Overview}
{\Large\bfseries\textcolor{CNavy}{Section 3.11: correspondence}}\\[3pt]
{\small\textbf{Classification:} Definition}\\
{\small\textbf{Learning objective:} Explain the precise meaning of correspondence and distinguish it from closely related architecture-description concepts.}
\end{CornellOverviewBox}
\vspace{3pt}

\begin{CornellNotesTable}
\CornellNoteRow{Term}{\textbf{correspondence}}
\CornellNoteRow{Definition}{An identified or named relationship among two or more architecture-description elements.}
\CornellNoteRow{Key characteristics}{\begin{itemize}[leftmargin=*,nosep,topsep=1pt]
\item For the purpose of correspondences, an architecture description (3.3) can be considered as an AD element (3.4) in another architecture description (3.3).
\item Correspondences can be identified or named relationship between architecture description elements (3.4) in different architecture descriptions (3.3) or between architecture description elements (3.4) stated in different notations.
\end{itemize}}
\CornellNoteRow{Distinction}{A correspondence is the relation itself; a correspondence method specifies rules or practices governing such relations.}
\CornellNoteRow{Example / application}{Correspondences are used to express a wide range of relationships, such as equivalence, composition, refinement, consistency, traceability, dependency, constraint, satisfaction, and obligation.}
\CornellNoteRow{Related sections}{3.4, 3.3}
\CornellNoteRow{Active recall}{\begin{itemize}[leftmargin=*,nosep,topsep=1pt]
\item How would you explain correspondence in your own words?
\item Which neighboring concept is easiest to confuse with correspondence, and how is it different?
\end{itemize}}
\end{CornellNotesTable}


\begin{CornellSummaryBox}{Summary}
An identified or named relationship among two or more architecture-description elements. A correspondence is the relation itself; a correspondence method specifies rules or practices governing such relations.
\end{CornellSummaryBox}

\end{document}
